% English translation of checked_pages/esq_119_012.tex; source: cote 119, site 12, PDF page 13, typescript page 10.
% AI-translated by GPT-6 Astra (OpenAI, ChatGPT Pro mode).
\scanpage{119/12}
\mspage{10}

of the isotopy classes of embeddings of such 1-complexes,
or of compact surfaces with boundary, in a compact oriented surface,
in terms of certain very simple combinatorial invariants and the
fundamental groups of the participants. This theorem, which should
extend without difficulty to embeddings of any compact space
(triangulable, to simplify matters) in a compact oriented surface,
recovers as easy corollaries several deep classical results
of surface theory, and in particular Baer's isotopy theorem.
It is finally going to be published, separately from the rest (and
ten years later, given the harshness of the times...), in Topology. Ladegaillerie's
work also contains a purely algebraic description,
in terms of fundamental groupoids, of the
``isotopy'' category of compact surfaces $X$, equipped with a topological
1-complex $K$ embedded in $X$. This description, which had the misfortune
to go against the ``taste of the day'' and consequently seems unpublishable,
has nevertheless served (and still serves) as a valuable guide in
my subsequent reflections, particularly in the context of absolute
algebraic geometry of characteristic zero.

The case in which $(X,K)$ is a 2-dimensional ``map'',\tnote{TN-012-01}{``Map'' here translates ``carte'', a surface equipped with the embedded graph and cell conditions specified in the sentence.}
i.e. in which the connected components of $X\backslash K$ are open 2-cells \unsure{.}
(and in which $K$ is moreover equipped with a finite set $S$ of ``vertices'',
such that the connected components of $K\backslash S$ are open 1-cells)
gradually attracts my attention in the following years.
The isotopy category of these maps admits a particularly
simple algebraic description, through the set of ``frames'' (or ``flags'',
or ``biarcs'') associated with the map, which is naturally equipped
with the structure of a set with a group of operators, under the group
\[ \underline{C}_2 = \langle \sigma_0,\sigma_1,\sigma_2 \,|\, \sigma_0^{\,2} = \sigma_1^{\,2} = \sigma_2^{\,2} = (\sigma_0\sigma_2)^2 = 1 \rangle \ , \]
which I call the (unoriented) \uline{cartographic group} of dimension 2.
It admits as a subgroup of index 2 the \uline{oriented cartographic group},
generated by products of even numbers of the generators, which can
also be described as
\[ \underline{C}_2^{+} = \langle \rho_s,\rho_f,\sigma \,|\, \rho_s\rho_f = \sigma,\ \sigma^2 = 1 \rangle \ , \]

\margin{12}

